\documentclass[handout,x11names]{beamer}
\usepackage{etex,beamerfoils}
\usepackage{array,pifont,amsmath,latexsym,epsfig,amsfonts,yfonts}
\usepackage{amstext,amsopn,amsxtra,amsgen,amsbsy,amscd,amssymb,dcolumn,graphicx,setspace,pgfplots,booktabs,natbib}
\usepackage[normalem]{ulem}
\usepackage[all]{xy}
\beamertemplatetheoremsnumbered
\setbeamertemplate{caption}[numbered]
\usetheme{Copenhagen}
\definecolor{NewBlue}{RGB}{4, 114, 155}
\definecolor{magenta}{cmyk}{0,1,0,0}
\usecolortheme[named=NewBlue]{structure}
%\usecolortheme[named=SkyBlue4]{structure}

\newcommand{\D}{\displaystyle}
\newcommand{\E}{\begin{eqnarray*}}
\newcommand{\F}{\end{eqnarray*}}
\newcommand{\EE}{\begin{eqnarray}}
\newcommand{\FF}{\end{eqnarray}}
\newcommand{\IZ}{\begin{itemize}}
\newcommand{\ZI}{\end{itemize}}
\newcommand{\EN}{\begin{enumerate}}
\newcommand{\NE}{\end{enumerate}}
\newcommand{\itemc}{\item[$\circ$]}
\newcommand{\IZdash}{\begin{itemize} \renewcommand{\labelitemi}{-}}
\newcommand{\tn}{\textnormal}

% New commands to assist Beamer
\newcommand{\itemi}{\item<1->}
\newcommand{\itemii}{\item<2->}
\newcommand{\itemiii}{\item<3->}
\newcommand{\itemiv}{\item<4->}
\newcommand{\itemv}{\item<5->}
\newcommand{\itemvi}{\item<6->}
\newcommand{\itemvii}{\item<7->}
\newcommand{\itemviii}{\item<8->}
\newcommand{\itemix}{\item<9->}
\newcommand{\itemx}{\item<10->}

\newcommand{\itemci}{\item[$\circ$]<1->}
\newcommand{\itemcii}{\item[$\circ$]<2->}
\newcommand{\itemciii}{\item[$\circ$]<3->}
\newcommand{\itemciv}{\item[$\circ$]<4->}
\newcommand{\itemcv}{\item[$\circ$]<5->}
\newcommand{\itemcvi}{\item[$\circ$]<6->}
\newcommand{\itemcvii}{\item[$\circ$]<7->}
\newcommand{\itemcviii}{\item[$\circ$]<8->}
\newcommand{\itemcix}{\item[$\circ$]<9->}
\newcommand{\itemcx}{\item[$\circ$]<10->}


\newcommand{\fr}{\frame}
\newcommand{\ft}{\frametitle}
\newcommand{\ons}{\onslide}
\newcommand{\tc}{\textcolor}

\pgfplotsset{height =7cm,compat=1.7,width=7cm,
    standard/.style={
    	unit vector ratio*=1 1 1,
        axis x line=middle,
        axis y line=middle,
        enlarge x limits=0.15,
        enlarge y limits=0.15,
        every axis x label/.style={at={(current axis.right of origin)},anchor=north west},
        every axis y label/.style={at={(current axis.above origin)},anchor=north east}
    }
}
\newcommand{\drawge}{-- (rel axis cs:1,0) -- (rel axis cs:1,1) -- (rel axis cs:0,1) \closedcycle}
\newcommand{\drawle}{-- (rel axis cs:1,1) -- (rel axis cs:1,0) -- (rel axis cs:0,0) \closedcycle}

\usetikzlibrary{patterns}
\usetikzlibrary{arrows}
\usetikzlibrary{shapes}

\renewcommand{\bibsection}{}

% Watermark example: https://texblog.org/2012/02/17/watermarks-draft-review-approved-confidential/
\usepackage{draftwatermark}
\SetWatermarkText{This content is protected and may not be shared uploaded or distributed.}
\SetWatermarkScale{0.25}
\SetWatermarkAngle{0}
\SetWatermarkColor[gray]{1.0}
\setbeamercolor{background canvas}{bg=}%transparent canvas

\title[Summers' Response\hspace{2em}\insertframenumber/
\inserttotalframenumber]{Lawrence Summers' Response to Prescott's Article}
\author{Brian C.~Jenkins \\ ~ \\University of California, Irvine}

\begin{document}

\frame{\maketitle}

\fr{
\ft{Introduction}

\IZ
\defcitealias{Prescott1986}{article by Edward Prescott} 
\defcitealias{Summers1986}{Lawrence Summers: ``Some Skeptical Observations on Real Business Cycle Theory''} 
\item The Fall 1986 issue of the Minneapolis Fed's \emph{Quarterly Review} featured an article by \citetalias{Prescott1986} summarizing the real business cycle (RBC) research program.

\

\itemii He argued that a stochastic neoclassical growth model could match empirical facts about the US business cycle.

\

\itemiii His framework has the appealing features that it is (relatively) simple and based on microeconomic principles.



\ZI
}



\fr{
\ft{Introduction}

\IZ
\item Prescott concludes on page 21 that:

\IZ
\itemii[] \emph{Economic fluctuations are optimal responses to uncertainty in the rate of technological change.}
\ZI

\onslide<3->{and that the:}

\IZ
\itemiv[] \emph{attention [of policymakers] should be focused not on fluctuations in output but rather on determinants of the average rate of technological advance.}
\ZI

\

\itemv Both points were and are controversial to many macroeconomists
\ZI
}


\fr{
\ft{Summers' Response}

\IZ
\defcitealias{Summers1986}{Summers' article, ``Some Skeptical Observations on Real Business Cycle Theory''} 
\item The issue of the \emph{Quarterly Review} in which Prescott's article appeared also featured a reaction by Lawrence Summers.

\

\itemii \citetalias{Summers1986}, presents sharp criticisms of RBC theory.
\ZI
}


\fr{
\ft{Summers' Response}

\IZ
\item In the opening of the article, Summers asserts that RBC theory is at odds with

\IZ
\itemii[] \emph{propositions thought self-evident by many macroeconomists and all of those involved in forecasting and controlling the economy on a day-to-day basis.}
\ZI

\

\itemii Summers provides four specific critiques of Prescott's framework
\ZI
}

\fr{
\ft{Summers' Response}

\IZ
\item Summary of Summers' critiques:

	\EN
	\itemii The parameterization that Prescott uses in the simulations are suspect.
	\itemiii Prescott doesn't provide supporting evidence for the existence of TFP shocks.
	\itemiv Prescott doesn't explore whether his RBC model can match price (e.g., real wages, interest rates) data.
	\itemv Prescott's model completely disregards the large-scale market failures that accompany business cycle downturns (e.g., unemployment)
	\NE
\ZI
}



\fr{
\ft{Summers' Response}

\IZ
\item Summers claims that progress in macroeconomics depends on

\IZ
\itemii[] \emph{the development of theories that can explain why exchange sometimes works well and other times breaks down.}
\ZI

\

\itemiii And he concludes with:

\IZ
\itemiv[] \emph{Nothing could be more counterproductive in this regard than a lengthy professional detour into the analysis of stochastic Robinson Crusoes.}
\ZI

\

\itemv In the first point he was prescient but in the last point he was greatly mistaken

\

\itemvi Contemporary business cycle models incorporate market failures into RBC-style models with great success.

\ZI


}



\fr{
\ft{References}
\bibliographystyle{aer}
\bibliography{comp_econ_references.bib}
}

\end{document}